% =====================================================================
%  19.5  Static and Offline Range Problems            sec:c13-static
%  Exports: prop:c13-offline  lst:c13-firstge  lst:c13-940-soln
%           prob:c13-cashback
%  Solved: CF 940E Cashback (2000), live-verified (statement + all 4 samples).
%    Task: array a[1..n], int c.  value of a run of length k = (sum of its
%    elements) - (the floor(k/c) SMALLEST elements).  Partition a into
%    contiguous runs to MINIMISE the total value.  1<=n,c<=1e5, 1<=a_i<=1e9.
%    Editorial structure (verified): WLOG every run has length 1 or c, so
%      dp[i] = min( dp[i-1] + a_i,                     // length-1 run
%                   dp[i-c] + (sum(a[i-c+1..i]) - min) )// length-c run, frees min
%    The window minimum is a STATIC range-min oracle -- the segment tree built
%    once over a (\coderef{c13-build}/\coderef{c13-query}), no updates.
%    Samples (verbatim from CF) and DP hand-checks:
%      1) n=3 c=5 [1 2 3]                       -> 6   (c>n: no discount, sum=6)
%      2) n=12 c=10 [1 1 10x6 9 10x3]           -> 92
%      3) n=7 c=2 [2 3 6 4 5 7 1]               -> 17  (traced below)
%      4) n=8 c=4 [1 3 4 5 5 3 4 1]             -> 23  (traced below)
%    DP for sample 3: dp=(0,2,3,8,9,13,16,17); partition [2 3][6 4][5 7][1].
%    DP for sample 4: dp=(0,1,4,8,12,15,18,22,23).  Both reproduce the answers.
%  Practices: CF 474F Ant colony (2100), CF 438D The Child and Sequence (2300),
%    both live-verified (title/rating/task shape/samples).
%  Running numbers agree with 00_opener: the worked static-descend example uses
%    19.1's max-tree over the shared 8-array (3,2,4,5,1,1,5,3).
% =====================================================================
\section{Static and Offline Range Problems}\label{sec:c13-static}

\begin{previouslybox}[Ch.~18 answered static ranks with a BIT; some folds need a tree]
Chapter~18 handled a \emph{static} order statistic -- ``how many values are
\(\le k\)'' -- with a frequency BIT, and Chapter~5's prefix sums answer a
\emph{static} range sum in \(O(1)\).  But a growing slice of range problems ask
for a fold that has \emph{no inverse and no monotone shortcut}: the minimum of a
window feeding a dynamic program, the leftmost position clearing a threshold,
the \(\gcd\) of a range counted against its members.  None of those telescope, so
the BIT of Ch.~18 cannot host them.  A segment tree \emph{built once} and then
queried many times is exactly the right object: an \(O(n)\) build buys an
\(O(\log n)\) oracle for \emph{any} associative fold, including the
threshold-descend of \(\secref{c13-search}\).  This section makes the segtree the
engine inside an offline sweep and inside a DP, and pays the debt Ch.~18 and
Ch.~5 left when their invertible tools ran out.
\end{previouslybox}

\begin{topicintro}
\term{Static} here means the array is fixed for the whole run -- the tree is
built once and never updated -- and \term{offline} means we may reorder or sort
the queries before answering them.  Both freedoms make the segment tree a pure
\emph{lookup oracle}: the queries are read-only folds or read-only descends, so
the \(4n\) build of \(\secref{c13-core}\) is the only mutable moment.  The payoff
is a family of problems that neither a prefix array (dies to any fold that is
not a sum) nor a BIT (dies to any non-invertible combine) can serve, yet that a
single pre-built tree answers in \(\bigO{\log n}\) each.

\textbf{Prerequisites.} \(\secref{c13-core}\) (build once, \(\coderef{c13-query}\)
fold).  \(\secref{c13-combine}\) (a node may carry a record such as
\((\gcd,\text{count})\)).  \(\secref{c13-search}\) (the threshold descend used as
the worked oracle here).  Ch.~5 (\(\chref{c24}\)) \(\tin/\tout\) flattening for
the tree pointer at the end.  Nothing else.

\textbf{After this section you can}
\begin{itemize}[nosep]
  \item turn a static range-min/max tree into a ``leftmost \(\ge x\) in
        \([l,r]\)'' oracle and read it offline; \hfill\lv{3}--\lv{4}
  \item recognise when a dynamic program's transition cost is a range fold and
        replace a naive re-scan with a single segtree query; \hfill\lv{4}--\lv{5}
  \item prove a structural ``optimal partition uses pieces of size \(1\) or \(c\)''
        lemma and code the resulting DP; \hfill\lv{5}
  \item explain why a sparse table beats the segtree for pure static RMQ, and
        where the segtree still wins (updates, descends, records). \hfill\lv{4}
\end{itemize}
\end{topicintro}

\begin{recognitionbox}[Pattern card -- ``fixed array, many range lookups, no inverse'']
\reckey{the input array never changes but \(q\) range questions do; each
question is a fold \(\min/\max/\gcd\) or a threshold ``first/last position in
\([l,r]\) satisfying an aggregate predicate''; a dynamic program whose
transition is \(\mathit{dp}[i]=\min_j \mathit{dp}[j]+\mathcal{F}(j{+}1,i)\) for a
range fold \(\mathcal{F}\); the phrase ``partition into contiguous parts'' with a
per-part cost that is a sum minus an extremum.}
\recuse{
  \item static array + read-only folds/descends -- build once
        (\(\coderef{c13-build}\)), query \(\bigO{\log n}\) each; this section;
  \item a DP whose range cost is \(\text{sum}-\min\) or \(\max\) of a window --
        let the tree return the extremum, keep running sums in a prefix array;
  \item offline threshold sweeps -- sort the queries, let one monotone tree pass
        answer them (\(\propref{c13-offline}\));
  \item a \emph{pure} static min/max with zero updates -- consider Ch.~21's sparse
        table for \(O(1)\) (\(\chref{c19}\)); the segtree is the choice the moment
        an update, a descend, or a record appears.
}
\recnotuse{
  \item the fold is a plain sum with point adds -- prefix array or BIT
        (\(\chref{c14}\)) is smaller and answers in \(O(1)\)/\(\bigO{\log n}\);
  \item whole intervals are \emph{modified} -- lazy propagation, Ch.~20
        (\(\chref{c15}\)); a static tree has no range-write path;
  \item ``\(k\)-th smallest in a static range \([l,r]\)'' with many queries and no
        updates -- a merge-sort or persistent tree, Ch.~21 (\(\chref{c19}\)).
}
One-line test: \emph{array fixed, questions are non-invertible range folds or
thresholds?  Build one tree, then answer every question as a read-only query or
descend.}
\end{recognitionbox}

% ---------------------------------------------------------------
\subsection{A Static Range-Min Oracle That Also Descends}
% ---------------------------------------------------------------

The simplest static use is \(\coderef{c13-query}\) as a lookup table: build once,
and every ``what is the minimum of \(A[l..r]\)'' is a fold over the canonical
cover.  The more interesting static use is the \emph{descend}
of \(\secref{c13-search}\) re-pointed at a stored range: ``the leftmost position
of \([l,r]\) whose value is at least \(x\)'' is answered by pruning on a
\(\max\)-aggregate, with no updates at all.

\begin{lstlisting}[caption={leftmost position in a closed range whose stored max clears a threshold -- a static descend}, label={lst:c13-firstge}]
// MAX tree over A[1..n].  Returns the smallest p in [L,R] with A[p] >= x,
// or -1 if no element of [L,R] reaches x.  A static, read-only search.
int first_ge(int v, int l, int r, int L, int R, long long x){
    if (l > R || r < L) return -1;          // node wholly outside [L,R]
    if (mx[v] < x)      return -1;          // no element here can qualify  (PRUNE)
    if (l == r)         return l;           // l is inside [L,R] and A[l] >= x
    int mid = (l + r) / 2;
    int p = first_ge(2*v,   l,   mid, L, R, x);   // try the LEFT half first
    if (p != -1) return p;                  // leftmost lives left
    return first_ge(2*v+1, mid+1, r, L, R, x);    // else the right half
}
// O(log n) once you hit the boundary; the mx[] prune kills whole subtrees.
\end{lstlisting}

\begin{observation}
\(\coderef{c13-firstge}\) is the \(\secref{c13-search}\) descend with a range
guard bolted on: the two \code{-1} base cases (outside \([L,R]\), and
\(\mathit{mx}<x\)) discard a whole subtree in \(O(1)\), and because we always
probe the left child first, the first leaf we accept \emph{is} the leftmost
qualifying position.  It reads no updates, so the tree is built once and reused
for all \(q\) queries -- the essence of the static setting.
\end{observation}

\begin{tracebox}[Static max-tree: leftmost position with value at least 5]
On \(\figref{c13-minbuild}\)'s tree shape but folding \(\max\) over
\(A=(3,2,4,5,1,1,5,3)\): \(t_4{=}3,t_5{=}5,t_6{=}1,t_7{=}5,\)
\(t_2{=}5,t_3{=}5,t_1{=}5\).  Query ``leftmost \(A[p]\ge 5\) in \([1,8]\)''.
Root \(t_1\): \(\mathit{mx}{=}5\not<5\), not a leaf, \(\mathit{mid}{=}4\); probe
left \(t_2[1..4]\): \(\mathit{mx}{=}5\), probe its left \(t_4[1..2]\):
\(\mathit{mx}{=}3<5\) \(\Rightarrow\) \code{-1}; probe its right \(t_5[3..4]\):
\(\mathit{mx}{=}5\), probe left \(t_{10}[3]\): \(4<5\Rightarrow\) \code{-1}; probe
right \(t_{11}[4]\): \(5\ge5\), leaf \(\Rightarrow\) return \(4\).  The left
subtree of the root succeeded, so we never touch \(t_3\).  Answer \(p{=}4\)
(\(A_4{=}5\), the first \(5\), cross-checking a linear scan).
\end{tracebox}

\begin{proposition}{Static folds plus offline ordering beat a per-query rescan}{c13-offline}
Let a batch of read-only range questions (folds or threshold descends) be issued
on a fixed array of length \(n\).  Building one segment tree in \(\bigO{n}\)
answers the whole batch in \(\bigO{q\log n}\), against \(\bigO{qn}\) for naive
rescans.  When the questions are further \emph{offline} -- reorderable, or
monotone in a threshold -- sorting them lets a single monotone pass (adding
positions as they qualify) reuse the same tree incrementally, so the total is
\(\bigO{(n+q)\log n}\) regardless of interleaving.  No inverse or monotonicity of
the \emph{data} is needed -- only that no query writes to the array.
\end{proposition}
\begin{proofbox}
Correctness of the fold half is \thmref{c13-merge} (canonical cover) and of the
descend half is \thmref{c13-descend} restricted to a read-only tree with a range
guard (the Observation above).  For the offline bound: process the
sorted questions in order of a non-decreasing threshold; each array position is
activated at most once, and between activations only \(\bigO{\log n}\) nodes
change, so \(n\) activations plus \(q\) descends cost \(\bigO{(n+q)\log n}\).
The tree is not even required to hold the thresholded values -- it may be a
plain \(0/1\) activity tree driven by the sweep -- which is exactly how a
``first \(\ge x\)'' batch collapses to one pass.
\end{proofbox}

\noindent\textbf{Why not a sparse table?}  For a \emph{pure} static \(\min/\max\)
\(\gcd\) lookup with zero updates, Ch.~21's sparse table answers in \(O(1)\) after
\(O(n\log n)\) build -- asymptotically better than \(\bigO{\log n}\) per query.
The segment tree stays the right tool when (i) a point update \emph{can} happen
between queries, (ii) the query is a \emph{descend} (find a boundary, not a fold
value) where the sparse table has no clean hook, or (iii) the node must carry a
\emph{record} (\(\secref{c13-combine}\)) whose idempotent-only sparse merge does
not apply.  940E below is case (iii)-adjacent: its oracle is a static range min,
but it is consumed by a DP, and pairing the tree with prefix sums keeps the code
shared with \(\secref{c13-core}\).

% ---------------------------------------------------------------
\subsection{Problem Lab: CF 940E Cashback}
% ---------------------------------------------------------------

\solvedproblem{c13-cashback}{Cashback}{CF 940E}
  {https://codeforces.com/problemset/problem/940/E}{CF 2000}
  {structural partition lemma (pieces of length 1 or c) + a static range-min oracle + prefix sums feeding a linear DP}

\begin{problemstatement}[CF 940E Cashback]
You are given an array \(a\) of length \(n\) and an integer \(c\).  The
\emph{value} of an array \(b\) of length \(k\) is the sum of its elements except
for the \(\lfloor k/c\rfloor\) smallest.  For example, the value of
\([3,1,6,5,2]\) with \(c=2\) is \(3+6+5=14\).  Among all possible partitions of
\(a\) into contiguous subarrays, output the smallest possible sum of the values
of these subarrays.

\textbf{Input.} \(n,c\) with \(1\le n,c\le 10^5\); then \(n\) values
\(1\le a_i\le 10^9\).

\textbf{Output.} One integer: the minimum total value of a contiguous partition.

\textbf{Examples.}
\begin{verbatim}
Input:          Output:
3 5             6
1 2 3

Input:                          Output:
12 10                           92
1 1 10 10 10 10 10 10 9 10 10 10

Input:            Output:       Input:              Output:
7 2               17            8 4                 23
2 3 6 4 5 7 1                 1 3 4 5 5 3 4 1
\end{verbatim}
\end{problemstatement}

\begin{attemptgate}[Attempt it yourself first -- \lv{5}, give yourself one sitting]
Total payment is \(\sum a\) minus the total \emph{discount}; a run of length \(k\)
discounts the \(\lfloor k/c\rfloor\) smallest of that run.  Minimising payment is
maximising discount, over all contiguous partitions -- an exponential space.
\peek{what does a run of length \(<c\), of length \(c\), and of length \(2c\) each
discount?}
\nudge{Length \(<c\): \(\lfloor k/c\rfloor=0\), no discount -- those elements are
paid in full, so split them into singletons freely.  Length \(c\): exactly one
discount (its minimum).  Length \(2c\): two discounts, but the two smallest of one
big run is \emph{never better} than the minima of two length-\(c\) runs.}
\fullhint{WLOG every run has length \(1\) or \(c\) (prove it, next).  Then
\(\mathit{dp}[i]=\min(\mathit{dp}[i{-}1]+a_i,\ \mathit{dp}[i{-}c]+(\text{sum}
(a_{i{-}c{+}1..i})-\min(a_{i{-}c{+}1..i})))\).  The running sums come from a
prefix array; the window minimum is a static range-min query on a tree built once
over \(a\) (\(\coderef{c13-build}\), \(\coderef{c13-query}\)).  \(\bigO{n\log n}\).
Full program below.}
\end{attemptgate}

\solstep{Understanding}
A run of length \(k\) is paid in full except for its \(\lfloor k/c\rfloor\)
smallest elements.  Cut the array into runs; the objective is the sum of the
runs' payments.  Because a payment is ``sum minus \(\lfloor k/c\rfloor\) minima,''
the \(\min\) of a window is the only non-additive piece -- and a window's minimum
is exactly a static range fold.  Once we know the shape of an optimal run
(length~1 or~\(c\)), the whole problem is a one-dimensional DP whose second
transition costs \(\text{windowSum}-\text{windowMin}\), served by a
\(\bigO{\log n}\) query.

\solstep{Key Observations}
\begin{observation}
A run with \(k<c\) discounts nothing; its cost is \(\sum\) over its elements no
matter where it is cut.  Replace every such run by singletons at zero change --
so short runs are ``free'' structure we can ignore.
\end{observation}
\begin{observation}
A run of length \(c\) discounts exactly one element, its minimum.  This is the
\emph{cheapest} way to buy one discount: any longer run that discounts a single
element (\(c\le k\le 2c{-}1\)) could instead keep the length-\(c\) sub-run
containing that minimum as its own run and singleton-ise the rest, at equal cost.
\end{observation}
\begin{observation}
A run of length \(\ge 2c\) discounts \(\ge2\) elements, but only the run's very
smallest.  Splitting it into consecutive length-\(c\) runs discounts one minimum
per block, and the sum of per-block minima is \(\ge\) the block-of-\(2c\) two
smallest -- so lengthening past \(c\) never helps.  Hence some optimal partition
uses only runs of length \(1\) and \(c\).
\end{observation}

\solstep{Think It Through With Me}
\begin{thinkbox}
\deadend{Optimise the discount by a global ``sort and pick minima''}
{The runs must be \emph{contiguous} and a discount is the run's own minimum, so a
greedy that sorts candidate length-\(c\) windows by their minimum and takes
non-overlapping ones is \emph{wrong}: a slightly worse early minimum can enable a
far better later one.  (It fails even small cases -- you need the DP.)}
\deadend{Range-sum DP with a naive re-scan for the window minimum}
{Reading each window's minimum by scanning \(c\) elements inside the DP is
\(O(nc)\), which at \(n=c=10^5\) is \(10^{10}\).  A static range-min tree returns
that minimum in \(\bigO{\log n}\) for a total \(\bigO{n\log n}\).}
\aha{Collapse the partition to runs of length \(1\) or \(c\); then
\(\mathit{dp}[i]=\min(\mathit{dp}[i{-}1]+a_i,\ \mathit{dp}[i{-}c]+\mathit{sum}-
\mathit{min})\) with a once-built min-tree supplying \(\mathit{min}\) and a
prefix array supplying \(\mathit{sum}\).}
\end{thinkbox}

\solstep{The Approach Ladder}
\begin{center}\renewcommand{\arraystretch}{1.25}
\begin{tabular}{@{}c p{0.22\linewidth} p{0.33\linewidth} l c@{}}
\toprule
\textbf{\#} & \textbf{Approach} & \textbf{Idea} & \textbf{Time} & \textbf{Verdict}\\
\midrule
1 & \lv{3}\ enumerate partitions & all cut sets, sort each run & \(\bigO{2^n n\log n}\) & \tle \\
2 & \lv{4}\ full DP, naive minima & every \((j,i)\) run, scan min & \(\bigO{n^2 c}\) & \tle \\
3 & \lv{4}\ DP, sparse-table minima & \(O(1)\) static RMQ, all \(j<i\) & \(\bigO{n^2}\) & too many states \\
4 & \lv{5}\ \textbf{\(1\)-or-\(c\) DP + segtree} & two transitions, \(\bigO{\log n}\) window min & \(\bigO{n\log n}\) & \ac \\
\bottomrule
\end{tabular}\end{center}
The structural lemma is what turns the \(O(n^2)\)-transition DP into a
two-transition DP -- the segtree then makes each transition cheap.

\solstep{Solution}
Build a static \(\min\)-tree over \(a[1..n]\) once; fill \(\mathit{pref}[i]=
\sum_{p\le i}a_p\).  Set \(\mathit{dp}[0]=0\).  For \(i=1..n\): take the
length-~\(1\) option \(\mathit{dp}[i{-}1]+a_i\), and if \(i\ge c\) the length-~\(c\)
option \(\mathit{dp}[i{-}c]+(\mathit{pref}[i]-\mathit{pref}[i{-}c])-
\mathit{queryMin}(i{-}c{+}1,i)\); \(\mathit{dp}[i]\) is the smaller.  Print
\(\mathit{dp}[n]\).  (\(c>n\Rightarrow\) no length-~\(c\) transition ever fires and
the answer is \(\sum a\), which the code yields for free.)

\solstep{Proof of Correctness}
\begin{proofbox}
\textbf{Lemma (shape).} Some optimal partition has every run of length \(1\) or
\(c\).  By Observations: a run of length \(k<c\) costs \(\sum\), so replace it
with singletons (equal cost); a run of length \(c\le k\le 2c{-}1\) discounts one
element -- keep a length-\(c\) sub-run containing that discounted minimum and
singleton the other \(k-c\) positions, and the discount (a minimum of a subset)
is unchanged, so cost is equal; a run of length \(\ge 2c\) splits into length-\(c\)
blocks whose minima sum to \(\ge\) the original \(\lfloor k/c\rfloor\) smallest
(a per-block minimum is at least as large as the correspondingly-ordered global
smallest it replaces), so splitting never increases cost.  Iterating reduces every
optimal partition to pieces of size \(1\) or \(c\).\\
\textbf{DP correctness.} \(\mathit{dp}[i]\) is the min cost for \(a[1..i]\).  The
last run of an optimal \(1\)-or-\(c\) partition of the prefix either has length
\(1\) (cost \(a_i\), prefix \(i{-}1\) solved by \(\mathit{dp}[i{-}1]\)) or length
\(c\) (cost \(\sum_{i-c+1}^{i}-\min\), prefix \(i{-}c\) by \(\mathit{dp}[i{-}c]\)).
Taking the min of exactly these two is exhaustive over the reduced family.
\(\mathit{queryMin}\) returns that window minimum by \thmref{c13-merge}, and
\(\mathit{pref}\) gives the window sum, so both transition costs are exact.
\textbf{Bounds.} \(n\) states, each \(O(1)\) queries of \(\bigO{\log n}\):
\(\bigO{n\log n}\) time, \(\bigO{n}\) memory.  \(\mathit{dp}\) values \(\le n\cdot
10^9=10^{14}\) fit \code{long long}.
\end{proofbox}

\solstep{Implementation}
\begin{lstlisting}[caption={CF 940E full solution: static range-min tree + prefix sums + the length-1-or-c DP}, label={lst:c13-940-soln}]
#include <bits/stdc++.h>
using namespace std;
const int MX = 100005;
const long long INF = (1LL<<62);
int n, c; long long a[MX];
long long pref[MX], dp[MX], mn[4*MX];

void build(int v, int l, int r){               // static min-tree, built once
    if (l == r) { mn[v] = a[l]; return; }
    int m = (l + r) / 2;
    build(2*v, l, m); build(2*v+1, m+1, r);
    mn[v] = min(mn[2*v], mn[2*v+1]);
}
long long queryMin(int v, int l, int r, int ql, int qr){   // fold [ql,qr]
    if (ql > qr) return INF;
    if (ql == l && qr == r) return mn[v];
    int m = (l + r) / 2;
    return min(queryMin(2*v,   l,   m, ql,           min(qr, m)),
               queryMin(2*v+1, m+1, r, max(ql, m+1), qr));
}
int main(){
    scanf("%d%d", &n, &c);
    for (int i=1;i<=n;i++){ scanf("%lld", &a[i]); pref[i]=pref[i-1]+a[i]; }
    build(1,1,n);
    dp[0]=0;
    for (int i=1;i<=n;i++){
        dp[i] = dp[i-1] + a[i];                          // last piece length 1
        if (i >= c){                                     // last piece length c
            long long s  = pref[i] - pref[i-c];
            long long m  = queryMin(1, 1, n, i-c+1, i);  // frees the window min
            dp[i] = min(dp[i], dp[i-c] + s - m);
        }
    }
    printf("%lld\n", dp[n]);
}
\end{lstlisting}

\solstep{Trace the Official Samples}
\begin{tracebox}[CF 940E sample 3: n=7, c=2, a = 2 3 6 4 5 7 1]
Prefix sums \(0,2,5,11,15,20,27,28\).  DP (window \(c{=}2\), min of the two):
\(i{=}1\): only length-1, \(\mathit{dp}[1]=2\).
\(i{=}2\): len1 \(\mathit{dp}[1]{+}3{=}5\); len2 \(\mathit{dp}[0]{+}(5{-}2){=}3\)
\(\Rightarrow\mathit{dp}[2]{=}3\).
\(i{=}3\): len1 \(2{+}6{=}9\)? \(\mathit{dp}[2]{+}6{=}9\); len2
\(\mathit{dp}[1]{+}(9{-}3){=}8\) \(\Rightarrow\mathit{dp}[3]{=}8\).
\(i{=}4\): len1 \(8{+}4{=}12\); len2 \(\mathit{dp}[2]{+}(10{-}4){=}9\)
\(\Rightarrow\mathit{dp}[4]{=}9\).
\(i{=}5\): len1 \(9{+}5{=}14\); len2 \(\mathit{dp}[3]{+}(9{-}4){=}13\)
\(\Rightarrow\mathit{dp}[5]{=}13\).
\(i{=}6\): len1 \(13{+}7{=}20\); len2 \(\mathit{dp}[4]{+}(12{-}5){=}16\)
\(\Rightarrow\mathit{dp}[6]{=}16\).
\(i{=}7\): len1 \(16{+}1{=}17\); len2 \(\mathit{dp}[5]{+}(8{-}1){=}20\)
\(\Rightarrow\mathit{dp}[7]{=}17\).  Output \texttt{17}.  The chosen pieces are
\([2,3][6,4][5,7][1]\) -- each length-~\(2\) run frees its minimum.
\end{tracebox}
\begin{tracebox}[CF 940E sample 4: n=8, c=4, a = 1 3 4 5 5 3 4 1]
Now a discount needs a length-\(4\) run.  The DP ends \(\mathit{dp}=(0,1,4,8,
12,15,18,22,23)\).  At \(i{=}5\) the length-~\(4\) window \([3,4,5,5]\) has sum
\(17\), min \(3\), cost \(\mathit{dp}[1]{+}14{=}15\) beats five singletons; the
run \([1]\) at the front and the trailing \([3,4,1]\) (a \(<c\) run, paid in full)
are singletons.  Output \texttt{23}.  Note this beats the tempting
``take two length-~\(4\) blocks'' partition \([1,3,4,5][5,3,4,1]\), whose discount
is \(1{+}1{=}2\) and payment \(26{-}2{=}24>23\) -- the DP's freedom to \emph{skip}
a poor minimum is exactly what the naive window-sorting greedy loses.
\end{tracebox}

\solstep{Complexity}
\begin{complexitybox}
Build \(\bigO{n}\); \(n\) DP states each doing one \(\bigO{\log n}\)
\(\mathit{queryMin}\) and \(O(1)\) prefix arithmetic \(\Rightarrow\)
\(\bigO{n\log n}\) time, \(\bigO{n}\) memory (\(4n\) tree + three \(\bigO{n}\)
\code{long long} arrays \(\approx 5.6\) MB).  At \(n=10^5\) this is \(\sim
1.7\cdot10^6\) operations -- comfortably inside \(2\) s.  A static sparse table
would drop the query to \(O(1)\) (\(\bigO{n}\) total) but cannot answer a
descend and is Ch.~21's tool; the segtree keeps this program on the same three
functions as \(\secref{c13-core}\).
\end{complexitybox}

\begin{hardeningbox}
\begin{itemize}[nosep,leftmargin=1.4em]
  \item \textbf{Prove the shape lemma before coding.}  The whole speed-up is
        ``only length \(1\) and \(c\) runs matter''; without it you write the
        \(O(n^2)\) DP over every \(j<i\).
  \item \textbf{Use \code{long long} everywhere.}  Costs reach \(n\cdot10^9=10^{14}\);
        the range-min identity must be \(>\!\) any value (\(\coderef{c13-940-soln}\)'s
        \code{INF}), and an \code{int} \(\mathit{dp}\) silently overflows to a wrong
        (often negative) minimum.
  \item \textbf{Window is \([i{-}c{+}1,\,i]\), inclusive.}  A half-open slip here
        reads \(c{+}1\) or \(c{-}1\) elements, so the discount counts the wrong
        smallest set and the DP degrades by one position everywhere downstream.
  \item \textbf{\(c>n\) and \(c=1\) boundaries.}  \(c>n\): no length-~\(c\) transition
        fires, answer \(=\sum a\).  \(c=1\): every run discounts its full length,
        answer \(0\).  The single loop handles both with the \(i\ge c\) guard; do
        not add special cases that then disagree with it.
\end{itemize}
\end{hardeningbox}

% ---------------------------------------------------------------
\subsection{A Light Pointer: Ranges on Trees}
% ---------------------------------------------------------------

A subtree of a fixed tree \emph{is} an interval once you stamp it with Chapter~5's
\(\tin/\tout\) (Ch.~5, \(\chref{c24}\), owns the flattening; we do not re-derive
it).  So ``aggregate over the subtree of \(v\)'' becomes ``fold over
\([\tin(v),\tout(v)]\)'' and drops straight onto \(\coderef{c13-query}\): put each
node's value at slot \(\tin(v)\), build one static tree, answer with a range fold.
What the segtree adds over Ch.~18's BIT is the \emph{combine}: a subtree \(\min\),
\(\max\), or \(\gcd\) is a non-invertible fold the BIT cannot maintain, so the tree
is the engine Ch.~5's ``subtree \(=\) interval'' was pointing at.  Two limits are
honest: if the queries \emph{modify} a whole subtree, you need lazy propagation
(Ch.~20, \(\chref{c15}\)); if the tree is dynamic (cut/link), you need
Euler-tour or top trees (Ch.~6) and heavy-light decomposition
(Ch.~22, \(\chref{c16}\)) to turn root-paths into segtree ranges.  Those engines
live in those chapters -- here the takeaway is only that a flattened subtree range
is a static range problem of exactly this section's kind.

% ---------------------------------------------------------------
\subsection{Common Pitfalls}
% ---------------------------------------------------------------

\begin{warningbox}[Trap 1: trusting a window-sorting greedy]
940E's discount is ``pick the largest minima,'' which \emph{looks} greedy, but the
runs are contiguous and one discount forces a whole window: taking a slightly
worse early minimum can block a much better one.  Sample~4 refutes the greedy
(\(24\) vs the true \(23\)).  When a partition's pieces are contiguous and
interact, reach for a DP -- the segtree makes it fast, not greedy.
\end{warningbox}

\begin{warningbox}[Trap 2: the wrong identity for a min oracle]
\(\coderef{c13-940-soln}\)'s \code{queryMin} returns \code{INF} on an empty
overlap so \(\min\) ignores it -- exactly the \(+\infty\) identity of
\(\secref{c13-core}\).  Returning \(0\) (or a stale array value) makes every
window look like it contains a \(0\), and the DP discounts a phantom cheapest
element on every transition.
\end{warningbox}

\begin{warningbox}[Trap 3: re-building the tree inside the DP]
The array is \emph{static}; build once (\(\coderef{c13-build}\)) and treat
\(\mathit{queryMin}\) as pure.  A rebuild per state is \(\bigO{n}\) inside
\(\bigO{n}\) states -- \(\bigO{n^2}\) for no reason.  The offline/static setting
is precisely the case where the tree is write-free after construction.
\end{warningbox}

\begin{warningbox}[Trap 4: using the segtree where a sparse table is right]
For a \emph{pure} static \(\min/\max/\gcd\) lookup table with zero updates and
zero descends, Ch.~21's sparse table answers in \(O(1)\).  Carrying a
\(\bigO{\log n}\) segtree there is correct but a factor slow and more code.  Pick
the segtree when an update, a threshold descend, or a multi-field record appears
-- and name the sparse table otherwise (\(\chref{c19}\)).
\end{warningbox}

% ---------------------------------------------------------------
\subsection{Practice Problems}
% ---------------------------------------------------------------

\begin{stuckbox}
\begin{itemize}[nosep,leftmargin=1.4em]
  \item Have you proved a \emph{structural} reduction (e.g. pieces of length \(1\)
        or \(c\)) before choosing the DP transitions?
  \item Is the tree built once and read many, or are you wrongly rebuilding it?
  \item What is the identity for your fold, and does an empty overlap return it?
  \item Is your range-min window the right closed interval (exactly \(c\) leaves)?
  \item Would a sparse table do this lookup in \(O(1)\)?  If so, why the tree?
\end{itemize}
\end{stuckbox}

\begin{practiceset}
  \practice{Ant colony}{CF 474F}
    {https://codeforces.com/problemset/problem/474/F}
    {CF 2100 L4}
    {Static array, \(t\) queries \([l,r]\) and \emph{no updates}.  ``Freed'' ants
     are those whose strength equals \(\gcd(s_l..s_r)\); answer \(=\) \((r{-}l{+}1)
     - \#\{\text{elements equal to the range gcd}\}\).  A textbook static fold with a
     \emph{record} node: store \((\gcd,\ \text{count of elements equal to that } \gcd)\)
     per node and merge as \(\secref{c13-combine}\) prescribes (equal child gcds add
     counts; otherwise the merged \(\gcd\) keeps only the children whose \(\gcd\) it
     divides).  Build once, one query per \([l,r]\) -- this is exactly the
     static-oracle pattern of \(\secref{c13-static}\), and a BIT cannot fold a
     \(\gcd\).}
  \practice{The Child and Sequence}{CF 438D}
    {https://codeforces.com/problemset/problem/438/D}
    {CF 2300 L5}
    {Maintain \(a[1..n]\) under three operations: range \(\sum\) query, range
     \emph{modulo} (\(a_i := a_i \bmod x\) for \(i\in[l,r]\)), and point set.
     Store \((\sum,\max)\) per node; a modulo by \(x\) leaves every element
     \(<x\) untouched, so \emph{descend only into nodes whose } \(\max\ge x\) and
     fix leaves on the way down (the prune is the \(\secref{c13-search}\) first-\(\ge
     x\) idea used to bound the work -- each element is reduced \(\bigO{\log}\)
     times).  It is \emph{not} lazy propagation: the range update is applied by
     pruning, the point set rebuilds \((\sum,\max)\) up the path.  A genuine
     range-\emph{assign}/add is Ch.~20 (\(\chref{c15}\)); this one gets its speed
     from the max-prune, and answers can exceed \(32\) bits.}
\end{practiceset}

\paragraph{Back-cites and forwards.}
\(\secref{c13-core}\) supplies the static build and \(\coderef{c13-query}\) the
oracle; \(\secref{c13-combine}\) the record nodes (here \((\gcd,\text{count})\),
\((\sum,\max)\)); \(\secref{c13-search}\) the threshold descend reused by the
worked oracle and CF~438D.  Chapter~18 (\(\chref{c14}\)) owns the invertible
prefix sums the static case shares with, and its \(\secref{c14-offline}\) owns the
offline-sweep habit \(\propref{c13-offline}\) extends.  Chapter~5 (\(\chref{c24}\))
owns \(\tin/\tout\).  Forward: pure static \(O(1)\) RMQ and the merge-sort /
persistent tree are Ch.~21 (\(\chref{c19}\)); range \emph{updates} are Ch.~20's
lazy (\(\chref{c15}\)); subtree/path engines are Ch.~22 (\(\chref{c16}\)).

% ---------------------------------------------------------------
\begin{recallbox}[Section 19.5 in 10 seconds]
Fixed array, non-invertible range questions \(\Rightarrow\) build one tree,
answer read-only: a \emph{fold} via \(\coderef{c13-query}\), or a
\emph{threshold} via \(\coderef{c13-firstge}\) (prune on \(\mathit{mx}<x\), take
the left leaf).  Offline order turns \(q\) descends into one monotone pass
(\(\propref{c13-offline}\)).  When a DP's transition cost is \(\text{sum}-\min\)
over a window, a static range-min tree makes it \(\bigO{\log n}\).  CF~940E:
prove runs reduce to length \(1\) or \(c\), then
\(\mathit{dp}[i]=\min(\mathit{dp}[i{-}1]{+}a_i,\ \mathit{dp}[i{-}c]{+}\mathit{sum}{-}\mathit{min})\).
Sparse table is \(O(1)\) for \emph{pure} static (\(\chref{c19}\)); the tree wins
with updates, descends, or records.
\selftest{
  \item State the 940E shape lemma (only lengths \(1\) and \(c\)) and the one
        sentence for why a length-\(\ge 2c\) run never helps.
  \item What makes the window-sorting greedy fail on sample 4 -- give the numbers?
  \item Why is \code{INF} the right empty-overlap return for \code{queryMin}, and
        what breaks if you return \(0\)?
  \item When would a sparse table replace this section's tree, and when would it
        \emph{not}?
  \item Which Ch.~20 feature would CF~438D need if operation 2 were range \emph{add}
        instead of range modulo?
}
\end{recallbox}
